% =====================================================================
%  FORMULACION MINIMA
%  Salida coronaria: resistor de Starling intramiocardico
%  Modelo acoplado LV 0D <-> coronaria 3D (CoupledLV0DCoronary3D)
%  Requiere: amsmath, amssymb, booktabs
% =====================================================================
\documentclass[11pt]{article}
\usepackage{amsmath,amssymb}
\usepackage{booktabs}
\usepackage[margin=2.3cm]{geometry}
\newcommand{\dd}[2]{\frac{\partial #1}{\partial #2}}
\title{Salida coronaria como resistor de Starling intramioc\'ardico:\\
       formulaci\'on m\'inima}
\date{}
\begin{document}\maketitle

\begin{abstract}
\noindent
Se sustituye la condici\'on de salida coronaria del modelo acoplado 0D--3D por
una formulaci\'on de un \'unico estado por salida. El diagn\'ostico de la
ejecuci\'on anterior localiza dos defectos estructurales, no de calibraci\'on:
la compliancia epic\'ardica $C_a$, introducida \'unicamente para estabilizar el
acoplamiento retrasado, forma con la inertancia del dominio 3D un resonador de
factor de calidad $2$ que explica tanto el rizado observado como el pico
sist\'olico esp\'ureo del flujo de salida; y el drenaje venoso hacia una
presi\'on auricular fija, mientras la compresi\'on empuja el nodo microvascular
hasta $p_{im}$, multiplica por $7{,}6$ la presi\'on motriz venosa en s\'istole.
La correcci\'on de ambos --- resistencia de salida impl\'icita en lugar de nodo
capacitivo, y garganta colapsable de Starling en lugar de drenaje a presi\'on
fija --- elimina la ley de tubo colapsable, las multiplicidades en paralelo, la
impedancia caracter\'istica y las siete constantes asociadas. El estado por
salida pasa de dos a uno, el n\'umero de constantes libres de $13$ a $8$, y
desaparecen los tres solvers no lineales anidados por evaluaci\'on de residuo.
El cierre WRMS se conserva \'integramente y se eval\'ua por tabla universal.
\end{abstract}

\section{Diagn\'ostico}
\label{sec:diag}

Las tres anomal\'ias reportadas --- oscilaciones, flujo de salida m\'aximo en
s\'istole y flujo distal desmedido --- se miden directamente sobre
\texttt{CoronaryArtery.csv} ($484$ pasos, $\Delta t=10^{-3}$~s). Conviene
separar el caudal de salida $Q$ en sus dos contribuciones, que el balance de
nodo hace accesibles sin instrumentar el c\'odigo:
\begin{equation}
  Q=\underbrace{C_a\,\dot p_{ca}}_{\text{carga del capacitor}}
   +\underbrace{q_a}_{\text{perfusi\'on}},
  \qquad
  q_a=\dot V+q_v .
  \label{eq:split}
\end{equation}

\begin{center}\small
\begin{tabular}{lrrr}
\toprule
 & S\'istole & Di\'astole & Cociente\\
\midrule
$Q$ (salida 3D)              & $0{,}92$ & $1{,}29$ & $0{,}71$\\
$q_a$ (perfusi\'on)          & $0{,}58$ & $1{,}44$ & $0{,}40$\\
$q_v$ (venoso)               & $7{,}69$ & $1{,}63$ & $4{,}71$\\
$C_a\dot p_{ca}$             & $0{,}33$ & $-0{,}16$ & ---\\
\bottomrule
\end{tabular}\\[2pt]
\footnotesize mL/s; s\'istole $t\in[0{,}22,0{,}42]$, di\'astole $t<0{,}16$.
\end{center}

\paragraph{(a) La fasicidad arterial ya era correcta.} La perfusi\'on $q_a$
tiene cociente sist\'ole/di\'astole $0{,}40$, que es el valor fisiol\'ogico. Lo
que no es correcto es lo que ve el dominio 3D: en el instante $t=0{,}229$~s el
caudal de salida vale $2{,}30$~mL/s, de los cuales $2{,}08$ son carga de $C_a$ y
solamente $0{,}22$ son perfusi\'on. \emph{El m\'aximo sist\'olico del caudal de
salida es en un $90\%$ carga capacitiva impulsada por $\dot p_{ao}$ durante el
ascenso a\'ortico}, no perfusi\'on. La inversi\'on de fase no est\'a en la
fisiolog\'ia del modelo sino en el elemento que se interpuso entre ella y la
malla.

\paragraph{(b) Las oscilaciones son la resonancia $L_{3D}$--$C_a$.} Con
$C_a=A L_p/(\rho c^2)$ y $L_{3D}=\rho\,\ell/A$ la inertancia del dominio, el
\'area se cancela y
\begin{equation}
  f_0=\frac{c}{2\pi\sqrt{L_p\,\ell}}\approx51\ \text{Hz},
  \qquad
  \zeta=\frac{Z_c}{2}\sqrt{\frac{C_a}{L_{3D}}}=0{,}25,
  \qquad
  Q_{\text{res}}=2 .
  \label{eq:res}
\end{equation}
El espectro del caudal de salida en la ventana diast\'olica presenta sus picos
en $12{,}5$--$37{,}5$~Hz, y un modelo unidimensional del acoplamiento reproduce
$58$ inversiones de signo con frecuencia dominante $63$~Hz. La impedancia
caracter\'istica amortigua ($\zeta=0{,}25$) pero no elimina: un resonador de
$Q_{\text{res}}=2$ sigue timbrando cuatro periodos. Obs\'ervese que $f_0$ y
$\zeta$ son \emph{independientes del radio de la salida}, de modo que el
fen\'omeno no se corrige refinando la geometr\'ia ni distribuyendo el caudal de
otro modo.

\paragraph{(c) El flujo venoso desmedido es la ausencia de garganta.} En
s\'istole $p_{im}$ alcanza $10\,349$~Pa. La constante de tiempo del
compartimento, $\tau=C R_v\approx0{,}7$~s con la compliancia emergente
$C\approx10^{-9}$~m$^3$/Pa, es larga frente a la s\'istole, de modo que $p_{tm}$
apenas var\'ia y \emph{el nodo microvascular sigue a $p_{im}$ casi uno a uno}:
$p_c$ pasa de $1\,910$ a $10\,509$~Pa. La presi\'on motriz venosa, que se toma
como $p_c-P_{RA}$, se multiplica por
\[
  \frac{10\,509-600}{1\,910-600}=7{,}6 ,
\]
mientras la modulaci\'on de conductancia que deber\'ia compensarla solo aporta
un factor $0{,}69$: la ley de tubo con $K_v=200$~Pa y $n=3/2$ reduce el radio
venular de $1{,}117$ a $0{,}908$ veces el no tensionado, un $9\%$. El producto,
$7{,}6\times0{,}69=5{,}2$, es el cociente $4{,}7$ medido. El lecho se vac\'ia a
$9$~mL/s y pierde $1{,}9$ de sus $7{,}0$~mL en un solo latido.

\paragraph{Conclusi\'on.} Ninguno de los tres es un problema de valores. El
primero y el segundo tienen la misma causa --- un capacitor colocado en la
salida por razones num\'ericas --- y el tercero es la falta de la \'unica
propiedad constitutiva que hace acotado el drenaje de un vaso colapsable
inmerso en un tejido a presi\'on: que la garganta se cierra a $p_{im}$.

\section{Principio rector}

\begin{quote}
\emph{Un vaso colapsable embebido en un tejido a presi\'on $p_{im}$ drena
contra $p_{im}$, no contra su reservorio distal. La presi\'on tisular no es una
fuente a\~nadida al balance: es la referencia de todo el compartimento.}
\end{quote}

\noindent
Es el resistor de Starling de Permutt y Riley y la cascada vascular de
Downey y Kirk (1975). Tomarlo en serio tiene una consecuencia que ahorra la
mitad del modelo: si la garganta se cierra a $p_{im}$, la presi\'on transmural
$p_{tm}=p_c-p_{im}$ es una variable de relajaci\'on estable con constante de
tiempo $CR_v$, acotada por debajo por $0$ y por arriba por $R_v\max q_a$. La
ley de tubo colapsable, la ley de volumen no negativo, las \'areas no
tensionadas y las multiplicidades en paralelo exist\'ian todas para acotar una
variable que ahora est\'a acotada por construcci\'on.

\section{Formulaci\'on}

\subsection{Topolog\'ia}

\[
\text{salida 3D}\;
\xrightarrow[\ \text{impl\'icita}\ ]{R_a}\;
\underbrace{\bigl(p_c,\;C\bigr)}_{\text{microvascular},\ p_{im}(t)}\;
\xrightarrow[\ \text{garganta a }p_{im}\ ]{R_v}\;
P_{RA}
\]

\noindent
Un solo nodo, una sola inc\'ognita. No hay nodo epic\'ardico, no hay capacitor
en la salida y no hay impedancia caracter\'istica.

\subsection{Ecuaciones}

Sea $p_{tm}=p_c-p_{im}$ la presi\'on transmural microvascular, \'unico estado
por salida. Con $p_{im}=\alpha\,p_{LV}(t)$,
\begin{align}
  \boxed{\;C\,\frac{d p_{tm}}{dt} = q_a-q_v\;}&
  \label{eq:balance}\\[2pt]
  \boxed{\;q_v = \begin{cases}
      \dfrac{p_{tm}+p_{im}-\max(p_{im},P_{RA})}{R_v\,\Phi_v}
        & \text{lumen abierto,}\\[8pt]
      0 & \text{garganta cerrada,}
    \end{cases}\;}&
  \label{eq:starling}\\[2pt]
  \boxed{\;p_{out} = p_{im}+p_{tm}+R_a\,\Phi_a\,Q,
  \qquad q_a\equiv Q\;}&
  \label{eq:pout}
\end{align}
$\Phi_a$ y $\Phi_v$ son los multiplicadores reol\'ogicos de la
secci\'on~\ref{sec:reo}, iguales a la unidad en el punto de calibraci\'on. El
volumen almacenado es $V=C\,p_{tm}$, y $V\ge0$ se cumple porque
\eqref{eq:starling} anula el drenaje en cuanto $p_{tm}$ alcanza el cierre.

\paragraph{Lectura de \eqref{eq:pout}.} La presi\'on tisular aparece
expl\'icitamente en la presi\'on que ve la malla. Ah\'i est\'a, en una l\'inea,
todo el impedimento sist\'olico: en s\'istole el dominio 3D descarga contra
$p_{im}+p_{tm}\approx11\,200$~Pa en lugar de contra $2\,000$~Pa, y la ca\'ida
motriz arteriolar cae de $7\,400$ a $900$~Pa. No hace falta ninguna modulaci\'on
de calibre para producir la dominancia diast\'olica.

\paragraph{Lectura de \eqref{eq:starling}.} En s\'istole $\max(p_{im},P_{RA})=
p_{im}$ y por tanto $q_v=p_{tm}/(R_v\Phi_v)$: la presi\'on tisular se cancela y
el drenaje queda gobernado por la \emph{sola} presi\'on transmural, que es lenta.
En di\'astole $p_{im}<P_{RA}$, la garganta se abre y se recupera el drenaje
ordinario $q_v=(p_c-P_{RA})/(R_v\Phi_v)$. El paso entre ambos reg\'imenes no es
un interruptor a\~nadido: es la definici\'on del resistor de Starling.

\paragraph{D\'onde va la compuerta.} Un resistor de Starling es una v\'alvula
antirretorno \emph{solamente mientras est\'a colapsado}, de modo que la
condici\'on se impone sobre el estado del lumen y no sobre el signo de la
presi\'on motriz:
\begin{equation}
  \text{lumen abierto}
  \iff
  p_{im}\le P_{RA}
  \ \ \text{(no hay compresi\'on)}
  \quad\text{o}\quad
  p_{tm}>0
  \ \ \text{(garganta distendida)} .
  \label{eq:gate}
\end{equation}
La distinci\'on importa: acotar $q_v\ge0$ siempre que $p_c<p_{drenaje}$
prohibir\'ia el reflujo venoso con el vaso completamente abierto, que es una
condici\'on real (presi\'on auricular derecha elevada, insuficiencia
tricusp\'idea). El reflujo por la rama abierta se autolimita y no necesita cota
inferior: $q_v<0$ da $C\dot p_{tm}>0$, que eleva $p_{tm}$ y cierra el
gradiente; el jacobiano sigue siendo $C/\Delta t+G_v>0$ en ambas ramas. En el
punto de operaci\'on de este modelo la compuerta no llega a activarse
($p_{tm}\in[444,1540]$~Pa en todo el barrido de $\alpha$), de modo que es una
salvaguarda y no una restricci\'on que est\'e conformando el resultado.

\subsection{Discretizaci\'on}

Euler impl\'icito y Newton escalar sobre $p_{tm}$:
\begin{equation}
  R(p_{tm})=\frac{C\bigl(p_{tm}-p_{tm}^{n}\bigr)}{\Delta t}-Q+q_v(p_{tm})=0,
  \qquad
  R'=\frac{C}{\Delta t}+\dd{q_v}{p_{tm}} .
  \label{eq:newton}
\end{equation}
Con $q_v$ mon\'otona creciente, $R'>0$ siempre y la iteraci\'on converge
globalmente desde cualquier iterado. Una sola inc\'ognita, sin b\'usqueda de
l\'inea, sin resolver el \'area del tubo dentro del residuo y sin invertir la
ley de flujo para obtener la ca\'ida del conducto.

\section{Acoplamiento 3D--0D}
\label{sec:acopl}

La \'unica raz\'on de existir de $C_a$ era la estabilidad: aplicada con retraso
de un paso, una resistencia $R$ produce el factor de amplificaci\'on
$|1-R\Delta t/L_{3D}|$, inestable en cuanto $R>2Z_{3D}$, y la resistencia
arteriolar por rama, $R_a\approx2{,}7\times10^{10}$, lo supera en un factor
$6$. La soluci\'on no es interponer un capacitor sino no retrasar el t\'ermino.
La contribuci\'on de $R_a$ a la tracci\'on de Neumann,
\begin{equation}
  \int_\Gamma R_a\,Q\,(v\cdot n)\,d\Gamma
  \;\simeq\;
  R_a\,A_\Gamma\int_\Gamma (u\cdot n)(v\cdot n)\,d\Gamma ,
  \label{eq:implicita}
\end{equation}
es exacta para perfil plano, sim\'etrica y semidefinida positiva para cualquier
perfil, y se ensambla con el mismo patr\'on que la impedancia normal de entrada
ya presente en el c\'odigo. Con ella el factor de amplificaci\'on pasa a
$1/(1+R\Delta t/L_{3D})<1$: \textbf{incondicionalmente estable}, sin
restricci\'on sobre $\Delta t$ ni sobre $R_a$.

\begin{center}\small
\begin{tabular}{lrrr}
\toprule
$R$ (Pa\,s/m$^3$) & $Z_{3D}$ & expl\'icito $|g|$ & impl\'icito $|g|$\\
\midrule
$6{,}7\times10^{8}$ (solo $Z_c$)   & $4{,}2\times10^{9}$ & $0{,}84$ & $0{,}86$\\
$4{,}5\times10^{9}$ (total)        & $4{,}2\times10^{9}$ & $0{,}07$ & $0{,}48$\\
$2{,}7\times10^{10}$ (por rama)    & $4{,}2\times10^{9}$ & $\mathbf{5{,}40}$ & $0{,}14$\\
\bottomrule
\end{tabular}\\[2pt]
\footnotesize $\Delta t=10^{-3}$~s. La tercera fila es la que hac\'ia divergir
el c\'alculo y motiv\'o $C_a$.
\end{center}

\noindent
Suprimido el capacitor, desaparece el resonador: el mismo modelo
unidimensional pasa de $58$ inversiones de signo y rizado RMS $1{,}5\times10^{-2}$
mL/s a $5$ inversiones y $5{,}3\times10^{-3}$, y el contenido restante es el del
propio forzamiento a\'ortico. $Z_c$ deja de ser necesario: su papel era
amortiguar $C_a$, y una impedancia de onda sin la compliancia que la genera no
tiene sentido en un dominio de pared r\'igida.

\section{Flujo retr\'ogrado}
\label{sec:back}

Nada en la formulaci\'on acota el signo de $q_a$. La reversi\'on sist\'olica
epic\'ardica es una \emph{predicci\'on} de \eqref{eq:pout}: ocurre cuando
$p_{im}+p_{tm}>p_{ao}$, esto es, cuando la presi\'on tisular supera a la
a\'ortica durante la contracci\'on isovolum\'etrica. El \'unico par\'ametro que
la gobierna es $\alpha$.

\begin{center}\small
\begin{tabular}{lrrrr}
\toprule
$\alpha$ & $q_a^{\min}$ (mL/s) & $t/T$ & $q_a$ sis/dia & $p_{tm}$ (Pa)\\
\midrule
$0{,}50$ & $+0{,}054$ & $0{,}209$ & $0{,}58$ & $1125$--$1540$\\
$0{,}65$ & $+0{,}006$ & $0{,}213$ & $0{,}33$ & $895$--$1423$\\
$0{,}80$ & $-0{,}032$ & $0{,}216$ & $0{,}10$ & $684$--$1342$\\
$1{,}00$ & $-0{,}108$ & $0{,}220$ & $-0{,}18$ & $444$--$1261$\\
\bottomrule
\end{tabular}\\[2pt]
\footnotesize Por salida, forzado con $p_{ao}(t)$ y $p_{LV}(t)$ del registro.
\end{center}

\noindent
El valor por defecto $\alpha=0{,}65$ cae justo en el umbral, lo que es
coherente con su lectura como transmisi\'on promediada sobre el espesor: en la
arteria epic\'ardica humana normal el reflujo sist\'olico es peque\~no o
ausente, y se hace prominente en el territorio subendoc\'ardico y en las
perforantes septales, que es lo que reproducen $\alpha=0{,}8$--$1{,}0$. El
instante del m\'inimo, $t/T\approx0{,}21$, es la contracci\'on isovolum\'etrica.

\paragraph{Sobre acotar el reflujo.} No se impone ninguna cota inferior.
Someti\'endola a las cuatro preguntas de la secci\'on~\ref{sec:crit}, una cota
del tipo $q_a\ge-q_{\max}$ falla las cuatro: no procede de una medida
independiente; su valor tendr\'ia que elegirse para que no interfiera con el
resultado; su ausencia es perfectamente defendible porque la reversi\'on est\'a
determinada por $\alpha$ y por la fase de $p_{ao}$ frente a $p_{LV}$; y no se
habr\'ia introducido antes de observar un valor que incomodase. Su efecto
pr\'actico ser\'ia enmascarar una divergencia en lugar de reportarla. El
t\'ermino de estabilizaci\'on de reflujo del contorno, en cambio, s\'i es
leg\'itimo y ya est\'a: es disipativo, se activa solo con $u\cdot n<0$ y su
magnitud, $\tfrac12\sigma\rho|u|\sim26$~Pa\,s/m a $5$~cm/s, est\'a cuatro
\'ordenes por debajo del coeficiente de resistencia impl\'icita
$R_a\Phi_a A\approx4{,}1\times10^{5}$~Pa\,s/m, de modo que amortigua la
reversi\'on sin prohibirla. El t\'ermino impl\'icito \eqref{eq:implicita} es
lineal y sim\'etrico y act\'ua igual sobre ambos signos: es una resistencia, no
un diodo.

\section{Reolog\'ia}
\label{sec:reo}

El cierre WRMS se conserva sin aproximaci\'on. Escribiendo el caudal de
Hagen--Poiseuille con una viscosidad aparente, $Q=\pi R^4\Delta p/(8\mu_{ap}L)$,
y comparando con $Q=\pi R^3 I(\tau_w)/\tau_w^3$ se obtiene
\begin{equation}
  \boxed{\;\mu_{ap}(\tau_w)=\frac{\tau_w^{4}}{4\,I(\tau_w)}\;},
  \qquad
  \dot\gamma_{\text{nom}}=\frac{4Q}{\pi R^{3}}=\frac{\tau_w}{\mu_{ap}},
  \qquad
  I(\tau_w)=\int_0^{\tau_w}\!\tau^2\dot\gamma(\tau)\,d\tau .
  \label{eq:muap}
\end{equation}
$\mu_{ap}$ es funci\'on \emph{universal} de $\tau_w$ para una reolog\'ia dada:
no depende de $R$, de $L$ ni de la rama. Se tabula una vez en
$\dot\gamma_{\text{nom}}$ y el elemento 0D solo interpola:
\begin{equation}
  \Phi_\bullet(q)=\frac{\mu_{ap}\bigl(\dot\gamma_\bullet^{0}\,|q|/q^{0}\bigr)}
                       {\mu_{ap}\bigl(\dot\gamma_\bullet^{0}\bigr)},
  \qquad \bullet\in\{a,v\}.
  \label{eq:phi}
\end{equation}
Con $241$ nodos log-espaciados en $\tau_w\in[10^{-6},10^{4}]$ el error m\'aximo
de interpolaci\'on log--log es $0{,}09\%$ y el medio $0{,}011\%$. Frente a la
evaluaci\'on anterior --- Newton para $\dot\gamma_w$ m\'as RK4 de $100$ pasos
por cada llamada, dentro de un Newton $2\times2$ con b\'usqueda de l\'inea, dos
\'arboles y seis salidas --- el coste por paso cae en tres \'ordenes de
magnitud sin perder una cifra.

\paragraph{D\'onde queda la se\~nal reol\'ogica.} En $R_v\Phi_v$. Como
$\dot\gamma_v^{0}\approx400$~s$^{-1}$ y el cizallamiento venoso recorre varias
d\'ecadas a lo largo del ciclo, $\Phi_v$ modula la resistencia de drenaje; por
\eqref{eq:balance} esto desplaza el punto de operaci\'on $p_{tm}$, que es
observable directo. La verificaci\'on 0D da, multiplicando la viscosidad
venosa:
\begin{center}\small
\begin{tabular}{lrrr}
\toprule
$\mu_v/\mu_v^{0}$ & $p_{tm}$ (Pa) & $Q$ (mL/min) & $q_a$ sis/dia\\
\midrule
$1{,}0$ & $831$--$1649$ & $102{,}3$ & $0{,}34$\\
$1{,}5$ & $1271$--$2112$ & $96{,}2$ & $0{,}30$\\
$2{,}0$ & $1662$--$2511$ & $90{,}9$ & $0{,}26$\\
\bottomrule
\end{tabular}
\end{center}
\noindent
La hiperviscosidad distiende el lecho y reduce el caudal medio. Es el efecto
que la formulaci\'on anterior no consegu\'ia mostrar, y aparece ahora en una
sola variable de estado.

\section{Calibraci\'on}

Se ejecuta una vez con $\mu_N$ y se congela. Con
$\Delta P=p_{ar}(0)-P_{RA}$ y $r_{p,i}$ medido sobre la malla:
\begin{equation}
  w_i=\frac{r_{p,i}^{3}}{\sum_j r_{p,j}^{3}},
  \qquad
  Q_i=Q_{\text{tot}}\,w_i,
  \qquad
  C_i=C_{\text{tot}}\,w_i,
  \label{eq:cal1}
\end{equation}
\begin{equation}
  R_{v,i}=\frac{f_v\,\Delta P}{Q_i},
  \qquad
  R_{a,i}=\frac{(1-f_v)\,\Delta P}{Q_i},
  \qquad
  p_{tm,i}(0)=P_{RA}+f_v\Delta P-\alpha\,p_{LV}(0).
  \label{eq:cal2}
\end{equation}
Tres divisiones por salida. Todo lo dem\'as --- radios arteriolares y
venulares, longitudes efectivas, multiplicidades, \'areas no tensionadas,
rigideces de pared, tiempos de tr\'ansito --- desaparece: eran par\'ametros de
una geometr\'ia expl\'icita que el modelo ya no necesita, porque
$R_a$, $R_v$ y $C$ son las cantidades que realmente entran en
\eqref{eq:balance}--\eqref{eq:pout}. La constante de tiempo de vaciado,
$\tau=C_iR_{v,i}$, se reporta como predicci\'on verificable.

\section{Constantes}

\begin{center}\small
\begin{tabular}{llll}
\toprule
S\'imbolo & C\'odigo & Valor & Origen\\
\midrule
$\mu_0,\mu_\infty,\lambda,n,a$ & \texttt{viscosity} & --- & reolog\'ia medida\\
$\mu_N$ & \texttt{newtonianCalibrationViscosity} & 3{,}5e-3 Pa\,s & referencia\\
\midrule
$Q_{\text{tot}}$ & \texttt{lcaTargetFlow} & 2{,}0e-6 m$^3$/s & 120 mL/min\\
$f_v$ & \texttt{venularPressureFraction} & 0{,}13 & reparto medido de carga\\
$C_{\text{tot}}$ & \texttt{coronaryComplianceTotal} & 3{,}2e-10 m$^3$/Pa & compliancia identificada\\
$\alpha$ & \texttt{intramyocardialFraction} & 0{,}65 & transmisi\'on tisular\\
$P_{RA}$ & \texttt{rightAtrialPressure} & 600 Pa & $\approx4{,}5$ mmHg\\
$\dot\gamma_a^{0}$ & \texttt{arteriolarShearRate} & 800 s$^{-1}$ & $4\bar v_a/r_a$\\
$\dot\gamma_v^{0}$ & \texttt{venularShearRate} & 400 s$^{-1}$ & $4\bar v_v/r_v$\\
\bottomrule
\end{tabular}
\end{center}

\noindent
Ocho constantes libres m\'as la reolog\'ia medida. Suprimidas respecto de la
formulaci\'on anterior: $c$ (\texttt{pulseWaveVelocity}),
$K_a$, $K_v$ (rigideces de pared), $m$, $n$ y las cotas de \'area de la ley de
tubo, $\bar v_a$, $\bar v_v$, $T_a$, $T_v$, $r_a$, $r_v$, $L_a$, $L_v$,
$N_a$, $N_v$, $A_u^a$, $A_u^v$, $Z_c$, $C_a$ y la geometr\'ia del conducto
epic\'ardico. Cantidades derivadas: $R_{a,i}$, $R_{v,i}$, $C_i$, $\tau_i$.

\section{Verificaci\'on}

Modelo 0D forzado con $p_{ao}(t)$ y $p_{LV}(t)$ del registro, $8$ ciclos,
$\Delta t=5\times10^{-5}$~s, m\'etricas sobre el \'ultimo ciclo.

\begin{center}\small
\begin{tabular}{lrrrrr}
\toprule
$C$ (m$^3$/Pa) & $\tau=CR_v$ (s) & $q_a$ sis/dia & $q_v$ pico/media & $p_{tm}$ (Pa) & $\Delta V$ (mL)\\
\midrule
$1{,}0\times10^{-10}$ & $0{,}07$ & $0{,}38$ & $1{,}28$ & $587$--$1799$ & $0{,}12$\\
$3{,}2\times10^{-10}$ & $0{,}22$ & $0{,}34$ & $1{,}33$ & $831$--$1649$ & $0{,}26$\\
$1{,}0\times10^{-9}$  & $0{,}68$ & $0{,}31$ & $1{,}25$ & $1126$--$1450$ & $0{,}32$\\
$3{,}0\times10^{-9}$  & $2{,}03$ & $0{,}30$ & $1{,}18$ & $1246$--$1358$ & $0{,}34$\\
\bottomrule
\end{tabular}
\end{center}

\begin{enumerate}
\item \textbf{Fasicidad arterial correcta y robusta.} $q_a$ sis/dia
  $=0{,}30$--$0{,}38$ sobre \emph{treinta veces} el rango de $C$. No es un
  ajuste: es una consecuencia estructural de \eqref{eq:pout}.
\item \textbf{Flujo venoso acotado.} Pico/media $\le1{,}33$ frente al $4{,}7$
  medido en la ejecuci\'on anterior, y sin desplazamiento neto de volumen.
\item \textbf{$p_{tm}>0$ en todo el ciclo}, de modo que $V=Cp_{tm}\ge0$ se
  cumple sin ley colapsable, sin penalizaci\'on y sin $p_0$.
\item \textbf{Volumen desplazado por latido} $0{,}12$--$0{,}34$~mL, frente a
  $1{,}9$~mL en la ejecuci\'on anterior. Es el orden fisiol\'ogico del
  desplazamiento sist\'olico coronario.
\item \textbf{Conservaci\'on exacta}: $\langle q_a\rangle=\langle q_v\rangle$ a
  precisi\'on de m\'aquina.
\end{enumerate}

\section{Criterio de correcci\'on frente a ajuste}
\label{sec:crit}

Se aplican las cuatro preguntas del documento anterior a los dos cambios.

\paragraph{Garganta de Starling \eqref{eq:starling}.} Procede de la mec\'anica
del tubo colapsable (Permutt--Riley; Downey \& Kirk 1975), no de un valor
medido a posteriori. El comportamiento es robusto: $q_a$ sis/dia var\'ia entre
$0{,}30$ y $0{,}38$ en treinta veces el rango de $C$. Su ausencia es
f\'isicamente indefendible: hace drenar un vaso comprimido contra un
reservorio que no ve. Y deb\'ia haberse impuesto desde el momento en que el
compartimento se declar\'o intramioc\'ardico.

\paragraph{Resistencia impl\'icita \eqref{eq:implicita}.} No introduce ninguna
constante: es el mismo $R_a$, ensamblado sin retraso. Es incondicionalmente
estable en todo el rango, no en una ventana. Y el t\'ermino omitido es el
t\'ermino diagonal exacto de la linealizaci\'on de la propia condici\'on de
contorno, luego omitirlo era la inconsistencia.

\paragraph{Contraste.} Ajustar $K_v$ para reducir el pico venoso no las
satisface: no hay medida independiente que lo fije en el valor requerido, el
comportamiento exige un valor de filo de navaja porque el cociente $7{,}6$ de
presi\'on motriz ha de compensarse en un $9\%$ de radio, y el mecanismo real es
el cierre de la garganta, no la reducci\'on gradual del calibre.

\section{Limitaciones}

\paragraph{Pico venoso sist\'olico.} La formulaci\'on entrega $q_v$ pico/media
$\approx1{,}3$. El pico sist\'olico que se mide en el seno coronario procede de
la descarga del compartimento venoso \emph{epic\'ardico}, que est\'a fuera del
miocardio y no forma parte del modelo: aqu\'i $q_v$ representa la salida
venular intramural, que fisiol\'ogicamente s\'i es casi plana. Recuperarlo
exigir\'ia un segundo nodo $(p_{ev},C_{ev})$ aguas abajo de la garganta y dos
constantes m\'as; se deja fuera deliberadamente.

\paragraph{Estratificaci\'on transmural.} $\alpha$ es una transmisi\'on
promediada sobre el espesor. La diferencia subendocardio/subepicardio, que es
donde reside la vulnerabilidad isqu\'emica, requerir\'ia salidas por capa.

\paragraph{Alcance del modelo reol\'ogico.} Los efectos no newtonianos ligados
a la naturaleza particulada de la sangre (F\aa hr\ae us--Lindqvist) no est\'an
contenidos en Carreau--Yasuda ni, por tanto, en $\mu_{ap}$.

\end{document}
